\chapter{Apollo Planning模块深度分析}
\label{chap:apollo_planning}

本章以Apollo 11.0源码为对象，梳理Planning模块的核心组件及其相互关系，内容包括模块总体架构、\texttt{ReferenceLine}平滑、路径-速度解耦两阶段优化、\texttt{PathBoundsDecider}相关的横向边界构建、\texttt{PathDecider}决策逻辑以及\texttt{SpeedBoundsDecider}的ST边界构建。源码层面的梳理为第\ref{chap:framework}章的问题识别提供事实依据。

\section{Planning模块总体架构}

\subsection{模块入口与数据流}

Planning模块的主入口函数为\texttt{PlanningComponent::Proc()}，每个规划周期由Cyber RT调度执行。从输入消息到输出轨迹的完整数据流如图\ref{fig:planning_flow}所示。整个处理过程划分为三个层次的函数调用。最外层的\texttt{PlanningComponent::Proc()}负责从Cyber RT消息通道读取五类输入数据并完成路由检查和有效性校验。中间层的\texttt{OnLanePlanning::RunOnce()}完成车辆状态更新、参考线缓存读取、轨迹拼接和Frame初始化等准备工作。参考线的生成与平滑并不在RunOnce的调用栈内，而是由\texttt{ReferenceLineProvider}托管的独立异步线程\texttt{GenerateThread}以50\,ms周期持续执行，RunOnce仅通过\texttt{GetReferenceLines()}读取其最新缓存，并通过\texttt{UpdateVehicleState()}和\texttt{UpdatePlanningCommand()}向该线程注入触发输入。核心执行层的\texttt{OnLanePlanning::Plan()}通过\texttt{ScenarioManager}选择场景，再由Stage驱动Task列表按顺序执行，最终产生\texttt{ADCTrajectory}轨迹输出。

\begin{figure}[htbp]
\centering
\adjustbox{max width=\textwidth}{% Planning 模块单周期数据流
\begin{tikzpicture}[
    scale=0.85, transform shape,
    >=Stealth,
    % 输入小盒
    inp/.style={rectangle, rounded corners=1pt, draw=deepblue,
                fill=inputyellow, minimum width=1.65cm, minimum height=0.6cm,
                font=\footnotesize, align=center, inner sep=2pt},
    % 层内子步骤
    step box/.style={rectangle, rounded corners=2pt, draw=deepblue,
                     fill=lightblue, minimum width=1.8cm, minimum height=0.6cm,
                     font=\footnotesize, align=center, inner sep=2pt},
    % 改进层子步骤
    step imp/.style={rectangle, rounded corners=2pt, draw=deeppink,
                     fill=improvedpink, minimum width=1.8cm, minimum height=0.6cm,
                     font=\footnotesize, align=center, inner sep=2pt},
    % 层背景
    layer bg/.style={rectangle, rounded corners=6pt, fill=bglight,
                     draw=deepblue!50, inner sep=8pt},
    % 改进层背景
    layer imp/.style={rectangle, rounded corners=6pt, fill=improvedpink!15,
                      draw=deeppink!60, inner sep=8pt},
    % 输出节点
    out box/.style={rectangle, rounded corners=3pt, draw=deeppink,
                    fill=outputpink, minimum width=4cm, minimum height=0.7cm,
                    font=\small\bfseries, align=center},
]

% ===================== 输入层 (y = 0) =====================
\node[inp] (i1) at (-4,   0) {Prediction\\[-1pt]Obstacles};
\node[inp] (i2) at (-2,   0) {Chassis};
\node[inp] (i3) at ( 0,   0) {Localization};
\node[inp] (i4) at ( 2,   0) {TrafficLight};
\node[inp] (i5) at ( 4,   0) {Planning\\[-1pt]Command};
\node[font=\small, text=deepblue, anchor=east] at (-5.2, 0) {Cyber RT 输入};

% 汇聚点
\coordinate (merge) at (0, -1.0);
\foreach \i in {i1,i2,i3,i4,i5}{
    \draw[flow] (\i.south) -- (merge);
}

% ===================== Layer 1: Proc() (y = -2.2) =====================
\node[step box] (p1) at (-2.8, -2.2) {CheckRerouting};
\node[step box] (p2) at ( 0,   -2.2) {构建 LocalView};
\node[step box] (p3) at ( 2.8, -2.2) {CheckInput};

\draw[flow] (p1) -- (p2);
\draw[flow] (p2) -- (p3);

% 隐藏锚点，使三层 fit 宽度一致
\coordinate (L1pad-l) at (-4.4, -2.2);
\coordinate (L1pad-r) at ( 4.4, -2.2);

\begin{scope}[on background layer]
    \node[layer bg, fit=(L1pad-l)(L1pad-r)(p1)(p2)(p3),
          label={[font=\small, text=deepblue, anchor=east]left:Proc()}] (L1) {};
\end{scope}

\draw[flow, line width=1.2pt] (merge) -- (L1.north);

% ===================== Layer 2: RunOnce() (y = -4.2) =====================
\node[step box] (r1) at (-3.4, -4.2) {VehicleState\\[-1pt]更新};
\node[step box] (r2) at (-1.1, -4.2) {参考线平滑};
\node[step box] (r3) at ( 1.2, -4.2) {轨迹拼接};
\node[step box] (r4) at ( 3.4, -4.2) {Frame\\[-1pt]初始化};

\draw[flow] (r1) -- (r2);
\draw[flow] (r2) -- (r3);
\draw[flow] (r3) -- (r4);

\begin{scope}[on background layer]
    \node[layer bg, fit=(r1)(r2)(r3)(r4),
          label={[font=\small, text=deepblue, anchor=east]left:RunOnce()}] (L2) {};
\end{scope}

\draw[flow, line width=1.2pt] (L1.south) -- (L2.north);

% ===================== Layer 3: Plan() (y = -6.2) =====================
\node[step imp] (s1) at (-3.4, -6.2) {Scenario\\[-1pt]Manager};
\node[step imp] (s2) at (-1.1, -6.2) {Scenario\\[-1pt]::Process};
\node[step imp] (s3) at ( 1.2, -6.2) {Stage\\[-1pt]::Process};
\node[step imp] (s4) at ( 3.4, -6.2) {Task\\[-1pt]::Execute};

\draw[flowimp] (s1) -- (s2);
\draw[flowimp] (s2) -- (s3);
\draw[flowimp] (s3) -- (s4);

\begin{scope}[on background layer]
    \node[layer imp, fit=(s1)(s2)(s3)(s4),
          label={[font=\small, text=deeppink, anchor=east]left:Plan()}] (L3) {};
\end{scope}

\node[font=\footnotesize, text=deeppink!80!black, anchor=west] at (L3.east) {~~核心执行层};

\draw[flowimp, line width=1.2pt] (L2.south) -- (L3.north);

% ===================== 输出 (y = -8.0) =====================
\node[out box] (output) at (0, -8.0) {ADCTrajectory};

\draw[flowimp, line width=1.2pt] (L3.south) -- (output.north);

\end{tikzpicture}
}
\caption{Planning模块数据流}
\label{fig:planning_flow}
\end{figure}

Task列表内部不进行回溯迭代，各Task按配置顺序单向执行，前序Task的输出作为后续Task的输入。这种流水线式的结构对每个Task的正确性要求较高。若前序Task中由\texttt{PathBoundsDeciderUtil}承载的路径边界构建逻辑出现不必要的保守决策，后续Task无法修正，只能在其输出之上继续处理。这一特征决定了横向边界的构造方式对最终轨迹的影响难以在下游环节被弥补。

\subsection{DependencyInjector依赖注入}

Apollo通过\texttt{DependencyInjector}类实现依赖注入，在\texttt{PlanningComponent::Init()}中创建并传递给各组件。表\ref{tab:dependency_injector}列出了其注入的主要依赖，包括规划上下文、历史轨迹和车辆状态等共享信息。

\begin{table}[htbp]
\centering
\caption{DependencyInjector注入的依赖组件}
\label{tab:dependency_injector}
\begin{tabular}{ll}
\toprule
依赖组件 & \textbf{功能说明} \\
\midrule
PlanningContext & 规划上下文，存储跨周期状态信息 \\
FrameHistory & Frame历史记录，用于轨迹拼接和平滑 \\
History & \texttt{ADCTrajectory}历史记录 \\
EgoInfo & 自车信息，含前方障碍物距离与路由信息 \\
VehicleStateProvider & 车辆状态，含位置、速度、加速度与航向 \\
LearningBasedData & 学习模型数据，供强化学习模式使用 \\
\bottomrule
\end{tabular}
\end{table}

\subsection{场景管理机制}

\texttt{ScenarioManager}在每个规划周期检查当前环境状态，通过调用各Scenario的\texttt{IsTransferable()}方法判断是否需要切换场景。图\ref{fig:scenario_hierarchy}展示了Scenario-Stage-Task的三级层次结构。Apollo支持19种场景类型，表\ref{tab:scenarios}列出了主要场景及其触发条件。

\begin{figure}[htbp]
\centering
\adjustbox{max width=\textwidth}{% Scenario-Stage-Task层次结构图
\begin{tikzpicture}[
    scale=0.85, transform shape,
    level 1/.style={sibling distance=5.5cm, level distance=2cm},
    level 2/.style={sibling distance=3cm, level distance=2cm},
    level 3/.style={sibling distance=2.2cm, level distance=2cm},
    scenar/.style={rectangle, draw=deepblue, fill=deepblue!60, text=white, font=\small\bfseries, rounded corners, minimum width=2.5cm, minimum height=0.8cm, align=center},
    stag/.style={rectangle, draw=deepblue, fill=skyblue, font=\small, rounded corners, minimum width=2.2cm, minimum height=0.8cm, align=center},
    tas/.style={rectangle, draw=deepblue, fill=lightblue, font=\footnotesize, rounded corners, minimum width=1.8cm, minimum height=0.7cm, align=center},
    mgr/.style={rectangle, draw=deeppink, fill=improvedpink, font=\small\bfseries, rounded corners, minimum width=3.5cm, minimum height=0.9cm, align=center}
]

  % ScenarioManager
  \node[mgr] (mgr) at (0, 1.5) {ScenarioManager};

  % Scenarios
  \node[scenar] (lf) at (-4, -0.5) {LANE\_FOLLOW};
  \node[scenar] (tl) at (0, -0.5) {TRAFFIC\_LIGHT\\PROTECTED};
  \node[scenar] (ss) at (4, -0.5) {STOP\_SIGN\\UNPROTECTED};
  \node[font=\small, gray] at (7, -0.5) {$\cdots$（共19种）};

  % Stages for LANE_FOLLOW
  \node[stag] (stage1) at (-4, -2.5) {LaneFollow\\Stage};

  % Tasks
  \node[tas] (t1) at (-7.8, -4.5) {LaneFollow\\Path};
  \node[tas] (t2) at (-5.3, -4.5) {Path\\Decider};
  \node[tas] (t3) at (-2.8, -4.5) {SpeedBounds\\Decider};
  \node[tas] (t4) at (-0.3, -4.5) {PiecewiseJerk\\Speed};
  \node[font=\footnotesize, gray] at (1.6, -4.5) {$\cdots$};

  % Stages for STOP_SIGN
  \node[stag] (ss1) at (3, -2.5) {Approach};
  \node[stag] (ss2) at (5.5, -2.5) {Creep};

  % 连线
  \draw[-Stealth, thick, deeppink] (mgr) -- (lf);
  \draw[-Stealth, thick, deeppink] (mgr) -- (tl);
  \draw[-Stealth, thick, deeppink] (mgr) -- (ss);

  \draw[-Stealth, thick, deepblue] (lf) -- (stage1);
  \draw[-Stealth, thick, deepblue] (ss) -- (ss1);
  \draw[-Stealth, thick, deepblue] (ss) -- (ss2);

  \draw[-Stealth, thick, deepblue!60] (stage1) -- (t1);
  \draw[-Stealth, thick, deepblue!60] (stage1) -- (t2);
  \draw[-Stealth, thick, deepblue!60] (stage1) -- (t3);
  \draw[-Stealth, thick, deepblue!60] (stage1) -- (t4);

  % 层次标注
  \node[font=\small\bfseries, deepblue, rotate=90] at (-9.8, -0.5) {Scenario层};
  \node[font=\small\bfseries, deepblue, rotate=90] at (-9.8, -2.5) {Stage层};
  \node[font=\small\bfseries, deepblue, rotate=90] at (-9.8, -4.5) {Task层};

  % 虚线分隔
  \draw[dashed, gray] (-9.3, -1.5) -- (7.5, -1.5);
  \draw[dashed, gray] (-9.3, -3.5) -- (7.5, -3.5);
\end{tikzpicture}
}
\caption{Scenario-Stage-Task三级管理架构}
\label{fig:scenario_hierarchy}
\end{figure}

图\ref{fig:planning_arch}以LaneFollow场景的默认流水线配置为基准，展开了Stage内部\texttt{task\_list\_}的真实执行顺序。路径侧由四个互斥的路径生成Task即\texttt{LaneChangePath}、\texttt{LaneFollowPath}、\texttt{LaneBorrowPath}、\texttt{FallbackPath}择一执行，再串接\texttt{PathDecider}与\texttt{RuleBasedStopDecider}；速度侧依次为\texttt{SpeedBoundsDecider}（PRIORI）、\texttt{PathTimeHeuristicOptimizer}、\texttt{SpeedDecider}、\texttt{SpeedBoundsDecider}（FINAL）和\texttt{PiecewiseJerkSpeedOptimizer}，其中\texttt{SpeedBoundsDecider}在管线中出现两次。历史版本中作为独立Task的\texttt{PathBoundsDecider}在Apollo 11.0中主要表现为工具类\texttt{PathBoundsDeciderUtil}，被上述四个路径生成Task内部共同调用以构建横向边界$[l^{-}(s), l^{+}(s)]$，本文改进点即定位于该边界构建逻辑，图中以星标与粉色高亮标识。外部输入侧，\texttt{ReferenceLineProvider}在独立异步线程中以50\,ms周期持续生成参考线，\texttt{TrajectoryStitcher}基于前一帧轨迹完成时间拼接，二者共同馈入Stage的输入端；最终由\texttt{ReferenceLineInfo::CombinePathAndSpeedProfile}将路径与速度合成为\texttt{ADCTrajectory}输出。

\begin{figure}[htbp]
\centering
\adjustbox{max width=\textwidth}{% Planning模块总体架构图
\begin{tikzpicture}[
    node distance=0.9cm,
    scale=1.0, transform shape,
    manager/.style={rectangle, draw=deepblue, minimum width=3.5cm, minimum height=1.2cm, fill=deepblue!60, font=\normalsize\bfseries, align=center, text width=3.3cm, text=white},
    task/.style={rectangle, draw=deepblue, minimum width=3.8cm, minimum height=1.2cm, fill=lightblue, font=\normalsize, align=center, text width=3.6cm},
    improved/.style={rectangle, draw=deeppink, minimum width=3.8cm, minimum height=1.2cm, fill=improvedpink, font=\normalsize, align=center, text width=3.6cm},
    final/.style={rectangle, draw=deepblue, minimum width=5cm, minimum height=1.2cm, fill=outputpink, font=\normalsize\bfseries, align=center, text width=4.8cm},
    arrow/.style={thick, -Stealth, deepblue}
]

% 标题
\node[font=\Large\bfseries] (title) at (0, 2) {Planning 模块总体架构};

% 管理层
\node[manager] (scenario) at (-5.2, 0) {Scenario\\Manager};
\node[manager] (stage) at (0, 0) {Stage\\Manager};
\node[manager] (task_runner) at (5.2, 0) {Task\\Runner};

% 管理层标签
\node[font=\small, below=0.15cm of scenario] {场景识别切换};
\node[font=\small, below=0.15cm of stage] {阶段状态管理};
\node[font=\small, below=0.15cm of task_runner] {任务执行调度};

% 箭头 - 管理层
\draw[arrow] (scenario) -- (stage);
\draw[arrow] (stage) -- (task_runner);

% 分隔线
\draw[thick, dashed, deepblue] (-8, -2) -- (8, -2);
\node[font=\large\bfseries, color=deepblue] at (0, -2.5) {核心Task模块};

% 左列 - 路径相关
\node[task] (ref_provider) at (-4, -4) {Reference Line\\Provider};
\node[task] (path_decider) at (-4, -6) {Path Decider};
\node[task] (path_optimizer) at (-4, -8) {Piecewise Jerk\\Path Optimizer};

% 右列 - 速度相关
\node[task] (path_bounds) at (4, -4) {Path Bounds\\Decider};
\node[improved] (speed_bounds) at (4, -6) {Speed Bounds\\Decider（改进）};
\node[task] (speed_optimizer) at (4, -8) {Piecewise Jerk\\Speed Optimizer};

% 改进标记
\node[font=\small\bfseries, color=deeppink] at (7.5, -6) {$\leftarrow$ 改进重点};

% 最终输出
\node[final] (stitcher) at (0, -10) {Trajectory Stitcher};

% 箭头 - Task流程
\draw[arrow] (ref_provider) -- (path_decider);
\draw[arrow] (path_decider) -- (path_optimizer);
\draw[arrow] (path_bounds) -- (speed_bounds);
\draw[arrow] (speed_bounds) -- (speed_optimizer);

% 汇聚箭头
\draw[arrow] (path_optimizer.south) -- ++(0, -0.3) -| ([xshift=-0.7cm]stitcher.north);
\draw[arrow] (speed_optimizer.south) -- ++(0, -0.3) -| ([xshift=0.7cm]stitcher.north);

% 外框
\begin{scope}[on background layer]
\node[rectangle, draw=deepblue, dashed, inner sep=0.6cm, fit=(title)(scenario)(stage)(task_runner)(ref_provider)(path_bounds)(path_optimizer)(speed_optimizer)(stitcher), fill=bglight] {};
\end{scope}

\end{tikzpicture}
}
\caption{LaneFollow场景下Stage的Task流水线与改进点定位}
\label{fig:planning_arch}
\end{figure}

\begin{table}[htbp]
\centering
\caption{Apollo主要场景类型及触发条件}
\label{tab:scenarios}
\small
\begin{tabularx}{\textwidth}{lX}
\toprule
场景类型 & \textbf{触发条件} \\
\midrule
LANE\_FOLLOW & 默认场景，正常车道跟随 \\
BARE\_INTERSECTION\_UNPROTECTED & 无保护路口，无交通灯与停车标志 \\
STOP\_SIGN\_UNPROTECTED & 遇到停车标志，距离在配置阈值内 \\
TRAFFIC\_LIGHT\_PROTECTED & 遇到交通灯且有路权保护 \\
TRAFFIC\_LIGHT\_UNPROTECTED\_LEFT\_TURN & 无保护左转遇到交通灯 \\
PULL\_OVER & 接近目的地且距离在配置范围内 \\
EMERGENCY\_STOP & 紧急停车 \\
PARK\_AND\_GO & 停车后重新启动 \\
VALET\_PARKING & 代客泊车 \\
\bottomrule
\end{tabularx}
\end{table}

\section{ReferenceLine平滑}

参考线的平滑质量决定了Frenet坐标系下障碍物SL投影的精度，进而影响路径边界的构建结果。本节从生成流程、数学建模和数据结构三个角度梳理\texttt{ReferenceLine}的构造过程。

\subsection{参考线生成流程}

ReferenceLine是Planning模块的规划基准线，由\texttt{ReferenceLineProvider}在\texttt{Init()}阶段经\texttt{cyber::Async}启动专用线程\texttt{GenerateThread}，以50\,ms周期反复调用\texttt{CreateReferenceLine()}产出最新参考线并写入内部缓存，供主规划线程的\texttt{RunOnce}通过\texttt{GetReferenceLines()}读取。生成流程主要步骤如下。

\begin{enumerate}
\item 通过\texttt{PncMapBase::GetRouteSegments()}从Routing结果中提取基于车辆当前位置的路由段。

\item 将路由段转换为\texttt{hdmap::Path}，再转换为\texttt{ReferenceLine}。此时的参考线直接来自高精地图的车道中心线，存在噪声和不连续性。

\item 调用\texttt{GetAnchorPoints()}，按\texttt{max\_constraint\_interval}间隔在参考线上均匀采样锚点，默认取0.25\,m。每个锚点包含位置、航向角和横向边界。

\item 以锚点为约束，通过离散点偏差最小化QP优化生成平滑的参考线。

\item 验证平滑后的参考线偏差不超过阈值，并根据车辆位置裁剪。
\end{enumerate}

图\ref{fig:ref_smooth}展示了参考线平滑的完整过程。

\begin{figure}[htbp]
\centering
\adjustbox{max width=\textwidth}{% 参考线平滑过程示意图
\begin{tikzpicture}[scale=0.85, transform shape]

  % ===== 上半：原始 vs 平滑 对比 =====

  % 原始车道中心线（明显锯齿/折线感）
  \draw[thick, gray, dashed]
    plot[smooth, tension=0.15] coordinates
      {(0,0.1) (0.6,0.35) (1.0,-0.1) (1.6,0.5) (2.0,0.15) (2.6,0.7) (3.2,0.3)
       (3.8,0.85) (4.2,0.5) (4.8,1.0) (5.4,0.6) (6.0,1.15) (6.6,0.8) (7.2,1.3)
       (7.8,0.95) (8.4,1.45) (9.0,1.2)};

  % 锚点（5m间距采样）
  \foreach \x/\y in {0/0.05, 1.5/0.2, 3.0/0.5, 4.5/0.75, 6.0/1.0, 7.5/1.15, 9.0/1.2} {
    \fill[deeppink] (\x, \y) circle (2.5pt);
    \draw[deeppink!50, thin] (\x, {\y-0.35}) -- (\x, {\y+0.35});
  }

  % 平滑后的参考线（光滑曲线）
  \draw[very thick, deepblue]
    plot[smooth, tension=0.6] coordinates
      {(0,0.05) (1.5,0.2) (3.0,0.5) (4.5,0.75) (6.0,1.0) (7.5,1.15) (9.0,1.2)};

  % 横向约束范围标注
  \draw[lightgreen!70!black, thin, {Stealth}-{Stealth}] (6.0, 0.55) -- (6.0, 1.45);
  \node[lightgreen!70!black, font=\tiny, left] at (5.95, 1.0) {$\pm\Delta l$};

  % 偏差标注
  \draw[dangerred, thin, {Stealth}-{Stealth}] (2.6, 0.42) -- (2.6, 0.7);
  \node[dangerred, font=\tiny, right] at (2.65, 0.56) {偏差};

  % ===== 图例（右上角） =====
  \begin{scope}[shift={(9.5, 1.6)}]
    \draw[thick, gray, dashed] (0, 0) -- (0.8, 0);
    \node[font=\scriptsize, anchor=west] at (1.0, 0) {原始中心线};

    \draw[very thick, deepblue] (0, -0.4) -- (0.8, -0.4);
    \node[font=\scriptsize, anchor=west] at (1.0, -0.4) {平滑参考线};

    \fill[deeppink] (0.4, -0.8) circle (2.5pt);
    \node[font=\scriptsize, anchor=west] at (1.0, -0.8) {锚点};

    \draw[lightgreen!70!black, thin, {Stealth}-{Stealth}] (0.2, -1.3) -- (0.6, -1.3);
    \node[font=\scriptsize, anchor=west] at (1.0, -1.3) {横向约束 $\pm\Delta l$};
  \end{scope}

  % ===== 下半：5步流程（左右与上半齐平：0 ~ 12） =====

  \node[rectangle, rounded corners=3pt, draw=deepblue, fill=inputyellow,
        minimum height=0.7cm, font=\footnotesize, align=center, text width=1.6cm]
    (s1) at (0.8, -2.2) {获取\\路由段};

  \node[rectangle, rounded corners=3pt, draw=deepblue, fill=lightblue,
        minimum height=0.7cm, font=\footnotesize, align=center, text width=1.6cm]
    (s2) at (3.4, -2.2) {生成原始\\参考线};

  \node[rectangle, rounded corners=3pt, draw=deepblue, fill=lightblue,
        minimum height=0.7cm, font=\footnotesize, align=center, text width=1.6cm]
    (s3) at (6.0, -2.2) {按5m采样\\锚点};

  \node[rectangle, rounded corners=3pt, draw=deeppink, fill=improvedpink,
        minimum height=0.7cm, font=\footnotesize, align=center, text width=1.6cm]
    (s4) at (8.6, -2.2) {QP样条\\平滑};

  \node[rectangle, rounded corners=3pt, draw=deepblue, fill=outputpink,
        minimum height=0.7cm, font=\footnotesize, align=center, text width=1.6cm]
    (s5) at (11.2, -2.2) {验证\\与裁剪};

  \draw[-Stealth, thick, deepblue] (s1) -- (s2);
  \draw[-Stealth, thick, deepblue] (s2) -- (s3);
  \draw[-Stealth, thick, deepblue] (s3) -- (s4);
  \draw[-Stealth, thick, deeppink] (s4) -- (s5);

\end{tikzpicture}
}
\caption{参考线平滑过程示意图}
\label{fig:ref_smooth}
\end{figure}

\subsection{离散点偏差平滑的数学建模}

Apollo早期版本采用分段五次多项式（Quintic Spline）的QP样条方法平滑参考线。自Apollo 7.0起，默认配置已切换至\texttt{DiscretePointsReferenceLineSmoother}，底层调用\texttt{FemPosDeviationSmoother}，以有限元离散点位置偏差最小化为目标。在Apollo 11.0默认规划配置中，QP样条方案已被关闭，离散点平滑配置为唯一生效方案。

该方法将平滑问题建模为离散点的二次规划。给定$n$个锚点$\mathbf{p}_i = (x_i, y_i)$，目标函数为

\begin{equation}
\begin{aligned}
\min_{\mathbf{p}} \; J ={}& w_{fem} \sum_{i=1}^{n-2} \|\mathbf{p}_{i-1} - 2\mathbf{p}_i + \mathbf{p}_{i+1}\|^2 \\
&+ w_{ref} \sum_{i=0}^{n-1} \|\mathbf{p}_i - \mathbf{p}_{i,ref}\|^2 + w_{len} \sum_{i=0}^{n-2} \|\mathbf{p}_{i+1} - \mathbf{p}_i\|^2
\end{aligned}
\label{eq:refline_obj}
\end{equation}
其中$w_{fem}$为有限元位置偏差权重，默认取$10^{10}$；$w_{ref}$为参考点偏差权重，默认取1.0；$w_{len}$为路径长度正则化权重，默认取1.0。$w_{fem}$远大于其余两项，表明Apollo优先保证相邻三点的共线性即曲率平滑，在此前提下兼顾与原始参考线的贴合度。

约束条件包括两类。
\begin{enumerate}
\item 对每个锚点$i$，约束其横向偏移不超过车道边界
\begin{equation}
\|\mathbf{p}_i - \mathbf{p}_{i,ref}\| \leq \Delta l_i
\end{equation}
其中$\Delta l_i$根据车道宽度计算，默认最大0.5\,m，最小0.1\,m。首末锚点的横向边界强制为零以固定端点。

\item 参考线首末两个锚点的横向偏移约束为零，防止端点偏离车道中心。
\end{enumerate}

\subsection{参考线数据结构}

平滑后的参考线由一系列\texttt{ReferencePoint}构成，每个点包含的字段如表\ref{tab:refpoint}所示。

\begin{table}[htbp]
\centering
\caption{ReferencePoint数据结构}
\label{tab:refpoint}
\begin{tabular}{lll}
\toprule
字段 & \textbf{单位} & \textbf{说明} \\
\midrule
$x, y$ & m & 全局笛卡尔坐标 \\
$heading$ & rad & 航向角 \\
$kappa$ & 1/m & 曲率 \\
$dkappa$ & 1/m$^2$ & 曲率变化率 \\
$s$ & m & 沿参考线的累积弧长 \\
$lane$ & — & 关联的车道信息 \\
$l$ & m & 相对车道中心线的横向偏移 \\
\bottomrule
\end{tabular}
\end{table}

\section{路径-速度解耦两阶段优化}

Apollo 2.0至5.0版本曾采用名为EM（Expectation-Maximization）Planner\cite{fan2018baidu}的独立规划器，以路径-速度解耦的两阶段优化为核心思想。在Apollo 11.0中，EM Planner已不再作为独立规划器存在，其路径-速度解耦的两阶段架构被融入了默认规划器\texttt{PublicRoadPlanner}的Scenario-Stage-Task流水线中。当前版本的四个规划器为NaviPlanner、\texttt{LatticePlanner}、RTKReplayPlanner和\texttt{PublicRoadPlanner}，其中\texttt{PublicRoadPlanner}为默认规划器，通过Task流水线继承了EM Planner的两阶段优化思想。

路径-速度解耦的核心做法是将三维时空轨迹规划降维为两个串联的二维规划子问题。第一阶段在忽略时间维度的前提下直接在SL空间上规划横向路径$l(s)$，第二阶段在已知$l(s)$的基础上沿参考线弧长规划速度曲线$s(t)$。图\ref{fig:path_speed_decouple}给出了该流水线的整体结构。参考线、感知障碍物、预测轨迹和车辆状态四种输入在阶段一汇入\texttt{PathBoundsDecider}相关边界构建逻辑与\texttt{PiecewiseJerkPathOptimizer}，产出离散路径点序列；阶段二将该路径作为已知条件建立ST图，由DP启发式搜索给出粗解作为热启动，再交由QP精细优化输出最终速度曲线；阶段三的输出与路径合成\texttt{ADCTrajectory}。这种降维的代价是路径阶段必须独立做出左右绕行、跟随、超车等决策，难以随动态障碍物的时间演化作出调整，这一限制与第\ref{chap:framework}章所分析的问题直接相关。

\begin{figure}[htbp]
\centering
\adjustbox{max width=\textwidth}{% 路径-速度解耦三阶段架构（左→右流水线）
\begin{tikzpicture}[
    scale=0.82, transform shape,
    >=Stealth,
    inp/.style={rectangle, rounded corners=1pt, draw=deepblue,
                fill=inputyellow, minimum width=1.6cm, minimum height=0.55cm,
                font=\footnotesize, align=center, inner sep=2pt},
    sl box/.style={rectangle, rounded corners=2pt, draw=deepblue,
                   fill=white, minimum width=2.8cm, minimum height=0.55cm,
                   font=\footnotesize, align=center, inner sep=2pt},
    st box/.style={rectangle, rounded corners=2pt, draw=deepblue,
                   fill=white, minimum width=2.8cm, minimum height=0.55cm,
                   font=\footnotesize, align=center, inner sep=2pt},
    sl bg/.style={rectangle, rounded corners=6pt, fill=lightblue!40,
                  draw=deepblue!60, inner sep=8pt},
    st bg/.style={rectangle, rounded corners=6pt, fill=skyblue!35,
                  draw=deepblue!60, inner sep=8pt},
    stglabel/.style={font=\footnotesize\bfseries, text=deepblue},
    out box/.style={rectangle, rounded corners=3pt, draw=deeppink,
                    fill=outputpink, minimum width=2.4cm, minimum height=0.7cm,
                    font=\small\bfseries, align=center},
]

% ==================== 左列：输入（与其他阶段间隔等价） ====================
\node[inp] (i1) at (0.9, 1.5)  {参考线};
\node[inp] (i2) at (0.9, 0.5)  {感知障碍物};
\node[inp] (i3) at (0.9, -0.5) {预测轨迹};
\node[inp] (i4) at (0.9, -1.5) {车辆状态};

% ==================== 阶段一：SL路径规划 (x=5) ====================
\node[sl box] (s1a) at (5, 0.85)  {PathBoundsDecider\\[-1pt]构建SL边界};
\node[sl box] (s1b) at (5, -0.85) {PiecewiseJerkPath-\\[-1pt]Optimizer\enspace QP求解};

\draw[flow] (s1a) -- (s1b);

\coordinate (s1pad-t) at (5, 1.5);
\coordinate (s1pad-b) at (5, -1.5);

\begin{scope}[on background layer]
    \node[sl bg, fit=(s1pad-t)(s1pad-b)(s1a)(s1b)] (S1) {};
\end{scope}

\node[stglabel, anchor=south] at (S1.north) {阶段一：路径规划};
\node[font=\scriptsize, text=deepblue!70, anchor=north] at (S1.south) {SL空间};


% ==================== 阶段二：DP速度搜索 (x=10) ====================
\node[st box] (s2a) at (10, 0.85)  {ST图网格构建};
\node[st box] (s2b) at (10, -0.85) {DP启发式搜索};

\draw[flow] (s2a) -- (s2b);

\coordinate (s2pad-t) at (10, 1.5);
\coordinate (s2pad-b) at (10, -1.5);

\begin{scope}[on background layer]
    \node[st bg, fit=(s2pad-t)(s2pad-b)(s2a)(s2b)] (S2) {};
\end{scope}

\node[stglabel, anchor=south] at (S2.north) {阶段二：速度搜索};
\node[font=\scriptsize, text=deepblue!70, anchor=north] at (S2.south) {ST空间};

% ==================== 阶段三：QP速度优化 (x=15) ====================
\node[st box] (s3a) at (15, 0.85) {PiecewiseJerkSpeed-\\[-1pt]Optimizer};
\node[st box] (s3b) at (15, -0.85){QP精细优化};

\draw[flow] (s3a) -- (s3b);

\coordinate (s3pad-t) at (15, 1.5);
\coordinate (s3pad-b) at (15, -1.5);

\begin{scope}[on background layer]
    \node[st bg, fit=(s3pad-t)(s3pad-b)(s3a)(s3b)] (S3) {};
\end{scope}

\node[stglabel, anchor=south] at (S3.north) {阶段三：速度优化};
\node[font=\scriptsize, text=deepblue!70, anchor=north] at (S3.south) {ST空间};

% ==================== 输出 (x=19.5) ====================
\node[out box] (output) at (19.5, 0) {ADCTrajectory};

% ==================== 输入 → 阶段一 ====================
\draw[flow] (i1.east) -- (S1.west);
\draw[flow] (i2.east) -- (S1.west);
\draw[flow] (i3.east) -- (S1.west);
\draw[flow] (i4.east) -- (S1.west);

% ==================== 阶段间箭头 ====================
\draw[flow] (S1.east) -- node[above, font=\scriptsize, text=deepblue] {$l(s)$ 路径}
                          node[below, font=\scriptsize, text=deeppink] {SL$\,\to\,$ST} (S2.west);

\draw[flow] (S2.east) -- node[above, font=\scriptsize, text=deepblue] {速度走廊} (S3.west);

\draw[flow] (S3.east) -- (output.west);

\end{tikzpicture}
}
\caption{路径-速度解耦三阶段架构}
\label{fig:path_speed_decouple}
\end{figure}

\subsection{路径规划阶段}

路径规划阶段通过一系列Path Task生成局部路径，典型的Task包括\texttt{LaneFollowPath}即车道跟随路径、\texttt{LaneChangePath}即换道路径、\texttt{LaneBorrowPath}即借道路径。路径优化采用与参考线平滑类似的QP方法，目标函数包含平滑性、参考线偏离和边界约束等项。

按Apollo 11.0默认\texttt{PublicRoadPlanner}在LaneFollow场景下的流水线配置，路径阶段的执行逻辑如下。默认由\texttt{LaneFollowPath}承担路径生成，该Task内部依次执行三个步骤。

DecidePathBounds()负责调用\texttt{PathBoundsDeciderUtil}工具类，沿参考线在每个离散纵向位置$s$上聚合车道几何、静态障碍物投影和借道决策，构建横向可行走廊$[l^-(s), l^+(s)]$作为优化问题的硬约束。

OptimizePath()以$l(s)$的分段jerk最小化为目标，在离散化网格上将平滑性、参考线偏离和边界约束组织为稀疏二次规划，调用OSQP求解得到横向位置、横向速度和横向加速度三元组序列，输出候选路径。

AssessPath()调用\texttt{PathAssessmentDeciderUtil}对候选路径进行有效性筛查，包含终点可达性、与静态障碍物的碰撞检测以及通行带是否被完全阻塞等判定。若\texttt{LaneFollowPath}评估失败，则按优先级依次尝试\texttt{LaneBorrowPath}和\texttt{FallbackPath}，每个Task内部同样执行上述三步。

最终由\texttt{PathDecider}在选定的路径上遍历所有障碍物，根据其与路径的相对横向距离和语义属性生成\texttt{NUDGE}、\texttt{IGNORE}等横向决策标签，后续速度规划阶段将据此构建ST边界。

上述流程最终输出\texttt{PathData}，其中封装了离散路径点序列及其Frenet坐标表达，作为速度规划阶段构建ST图和跟踪参考的唯一几何输入。

\subsection{速度规划阶段}

速度规划分为粗优化和精优化两步。

粗优化阶段采用动态规划（Dynamic Programming, DP）搜索，在ST图的离散栅格上执行。栅格的时间维度按固定步长$\Delta t$均匀采样，空间维度分为dense和sparse两部分以平衡精度和效率。DP的代价函数包括

\begin{equation}
C_{total} = C_{obstacle} + C_{spatial} + C_{speed} + C_{accel} + C_{jerk}
\label{eq:dp_cost}
\end{equation}

其中各项代价的计算方式如下。

障碍物代价为
\begin{equation}
C_{obstacle} = \begin{cases}
+\infty, & \rho < 0, \\
w_{obs} \cdot c_{default} \cdot (d_{follow} - d_{actual})^2, & 0 \le \rho \le \varepsilon,
\end{cases}
\end{equation}

其中$\rho$为至障碍物边界的符号横向间距（障碍物内部$\rho<0$），$\varepsilon$为近边界阈值。

速度代价为
\begin{equation}
\begin{aligned}
C_{speed} = {}& w_{exceed}(v - v_{limit})^2 \\
&+ w_{low}|v - v_{cruise}| + w_{ref}|v - v_{ref}|
\end{aligned}
\end{equation}

加加速度代价为
\begin{equation}
j = \frac{s_4 - 3s_3 + 3s_2 - s_1}{(\Delta t)^3}
\end{equation}

DP通过递推计算每个栅格点的最小累积代价，并记录前驱节点，最终通过回溯获得粗略的最优速度曲线。

精优化阶段采用PiecewiseJerk QP，以DP结果为参考轨迹建立QP问题进行精细化优化
\begin{equation}
\begin{aligned}
\min \sum_{k=0}^{K-1} \Big[ &w_s(s_k - s_{ref,k})^2 + w_v(v_k - v_{ref,k})^2 \\
&+ w_a a_k^2 + w_j \left(\frac{a_{k+1} - a_k}{\Delta t}\right)^2 \Big]
\end{aligned}
\label{eq:piecewise_jerk}
\end{equation}

DP输出的速度曲线作为$v_{ref,k}$传入QP，为非凸的ST边界约束提供了暖启动，引导QP在正确的可行区域内求解。

\subsection{两阶段信息传递}

两阶段之间的信息传递是单向的，路径阶段输出\texttt{PathData}，包含离散路径点和Frenet坐标系路径，速度阶段在此路径的基础上构建ST图并进行速度优化。在Apollo的Task流水线中，各Task按配置顺序执行，不进行回溯迭代。LANE\_FOLLOW场景的完整Task流水线配置如代码\ref{lst:lane_follow_pipeline}所示。

\begin{lstlisting}[language={json},caption={LANE\_FOLLOW场景Task流水线配置},label={lst:lane_follow_pipeline}]
// scenarios/lane_follow/conf/pipeline.pb.txt
stage {
  name: "LANE_FOLLOW_STAGE"
  type: "LaneFollowStage"
  task { name: "LANE_CHANGE_PATH"            type: "LaneChangePath" }
  task { name: "LANE_FOLLOW_PATH"            type: "LaneFollowPath" }
  task { name: "LANE_BORROW_PATH"            type: "LaneBorrowPath" }
  task { name: "FALLBACK_PATH"               type: "FallbackPath" }
  task { name: "PATH_DECIDER"                type: "PathDecider" }
  task { name: "RULE_BASED_STOP_DECIDER"     type: "RuleBasedStopDecider" }
  task { name: "SPEED_BOUNDS_PRIORI_DECIDER" type: "SpeedBoundsDecider" }
  task { name: "SPEED_HEURISTIC_OPTIMIZER"   type: "PathTimeHeuristicOptimizer" }
  task { name: "SPEED_DECIDER"               type: "SpeedDecider" }
  task { name: "SPEED_BOUNDS_FINAL_DECIDER"  type: "SpeedBoundsDecider" }
  task { name: "PIECEWISE_JERK_SPEED"        type: "PiecewiseJerkSpeedOptimizer" }
}
\end{lstlisting}

其中，路径阶段包含四种路径生成Task，即换道、跟车道、借道、降级，以及路径决策和规则停车决策两个决策Task；速度阶段包含先验速度边界构建、DP启发式搜索、速度决策、最终速度边界构建和QP速度优化五个Task。\texttt{SpeedBoundsDecider}被调用两次，分别在DP搜索前后执行先验和最终的边界构建。

除\texttt{PublicRoadPlanner}外，Apollo另提供\texttt{LatticePlanner}\cite{werling2010}作为备选规划器，采用Frenet坐标系下的采样方法生成候选轨迹。本文研究聚焦于LANE\_FOLLOW场景下的路径边界决策改进，\texttt{LatticePlanner}不在讨论范围内。

\section{PathDecider决策分析}
\label{sec:path_decider}

\texttt{PathDecider}是路径阶段末端的决策节点，负责针对参考线上的每个障碍物给出横向或纵向的语义标签，供后续速度规划阶段构建ST边界时使用。本节先梳理其决策类型和优先级，再分析判定逻辑，最后专门讨论其中的路径边界构建部分。

\subsection{决策类型与优先级}

\texttt{PathDecider}的决策分为横向和纵向两类，各类型及其语义优先级如表\ref{tab:path_decisions}所示。

\begin{table}[htbp]
\centering
\caption{PathDecider决策类型}
\label{tab:path_decisions}
\begin{tabular}{lll}
\toprule
方向 & \textbf{决策类型} & \textbf{说明} \\
\midrule
\multirow{2}{*}{横向} & NUDGE（LEFT/RIGHT） & 横向避让，左绕或右绕 \\
& IGNORE & 横向安全，忽略 \\
\midrule
\multirow{5}{*}{纵向} & STOP & 停车让行，最高优先 \\
& YIELD & 减速让行 \\
& FOLLOW & 跟车 \\
& OVERTAKE & 超车 \\
& IGNORE & 纵向安全，忽略，最低优先 \\
\bottomrule
\end{tabular}
\end{table}

注，Apollo源码中决策类型定义为枚举类型，不使用显式数值优先级。实际处理中，STOP决策的优先级最高，一旦判定需要停车便会覆盖其他决策，IGNORE的优先级最低。

\subsection{决策逻辑}

\texttt{PathDecider}的决策逻辑基于障碍物的SL边界与当前路径的横向关系进行判断。

\begin{enumerate}
\item 若障碍物不在路径$s$范围内，即$s_{end} < s_{path\_start}$或$s_{start} > s_{path\_end}$，或横向距离足够远，即$|l_{obs} - l_{ego}| > w_{ego}/2 + d_{ignore}$，则判定为IGNORE。其中$d_{ignore} = 3.0$\,m。

\item 若障碍物是阻塞障碍物，或横向重叠区域超过阈值，即$l_{obs\_end} \geq l_{ego} - l_{min\_nudge}$且$l_{obs\_start} \leq l_{ego} + l_{min\_nudge}$，则判定为STOP。其中$l_{min\_nudge} = w_{ego}/2 + d_{static\_buffer}/2$，$d_{static\_buffer} = 0.3$\,m。

\item 若障碍物在可避让范围内，根据其相对于自车的横向位置决定LEFT\_NUDGE或RIGHT\_NUDGE。
\end{enumerate}

关键特征在于\texttt{PathDecider}仅处理静态障碍物，动态障碍物的处理交由速度规划阶段的ST图来完成。

\subsection{PathBoundsDecider相关边界构建}
\label{sec:pathbounds_decider}

在Apollo 11.0中，路径边界构建逻辑主要由\texttt{PathBoundsDeciderUtil}承载。为便于与历史版本和常用称呼保持一致，本文以下仍以\texttt{PathBoundsDecider}简称该类逻辑。其核心职责是构建路径的横向边界，合并策略如下。

对每个$s$位置，遍历所有障碍物，根据nudge方向更新路径边界

\begin{equation}                                        
  \begin{cases}
  l_{\mathrm{upper}} = \min(l_{\mathrm{upper}},\; l_{\mathrm{obs\_right\_bound}}), & \text{RIGHT\_NUDGE} \\                                                                                          
  l_{\mathrm{lower}} = \max(l_{\mathrm{lower}},\; l_{\mathrm{obs\_left\_bound}}), & \text{LEFT\_NUDGE}                                                                                               
  \end{cases}                                                                                                                                                                                        
\end{equation} 

这种取最小/最大的合并策略是保守的，多个障碍物的约束取交集，确保路径不与任何障碍物冲突，但可能导致可行域被严重压缩。

该边界构建逻辑的完整处理流程如下。
\begin{enumerate}
\item 初始化路径边界，将上下界分别设为道路左右边界$[l_{road}^-(s), l_{road}^+(s)]$；
\item 遍历参考线上的纵向采样点$s_i$；
\item 对每个$s_i$，调用\texttt{UpdatePathBoundaryBySLPolygon}函数，依次处理该截面上的各个障碍物；
\item 对每个障碍物，通过\texttt{SetNudgeInfo}函数判定避让方向。源码中该判定实际包含三层分支，近距离障碍物即纵向距离小于5\,m且横向接近参考线中心时以自车初始SL状态为参考，中距离障碍物即位于自车纵向范围内时以自车尾部横向位置为参考，远距离障碍物才以上一次nudge记录的横向位置为参考。本文关注的是远距离障碍物的兜底逻辑，其判据为障碍物中心与参考值的相对位置，若障碍物中心在参考线右侧则左绕即将右边界收缩，若在左侧则右绕即将左边界收缩；
\item 逐障碍物更新边界后，若出现$l^+(s_i) < l^-(s_i)$则标记为BLOCKED。
\end{enumerate}

图\ref{fig:path_bounds_flow}以流程图形式给出了该过程，可以看到避让方向判定和边界更新两个节点在循环体内被逐障碍物触发。图中BLOCKED检查节点位于每次更新之后，一旦触发将直接终止规划并向上层返回降级请求；若通过检查则继续处理下一个障碍物，直到遍历完毕输出完整的path\_boundary。图右侧标注同时提示了本流程的主要设计局限，即决策一旦在循环中做出便不可修正，后续的障碍物无论与已做决策是否存在空间冲突，均无法回头重新选择更合适的避让方向。

\begin{figure}[htbp]
\centering
\adjustbox{max width=0.8\textwidth}{% PathBoundsDecider 构建流程图
\begin{tikzpicture}[
    scale=1.0, transform shape,
    box/.style={rectangle, rounded corners=3pt, draw=deepblue, fill=lightblue, minimum width=4.2cm, minimum height=0.9cm, font=\small, align=center},
    boximp/.style={rectangle, rounded corners=3pt, draw=deeppink, fill=improvedpink, minimum width=4.2cm, minimum height=0.9cm, font=\small, align=center},
    dec/.style={diamond, aspect=2, draw=deepblue, fill=skyblue, minimum width=4.2cm, minimum height=0.9cm, font=\small, align=center, inner sep=1pt},
    term/.style={rectangle, rounded corners=8pt, draw=deeppink, fill=outputpink, minimum width=3.2cm, minimum height=0.8cm, font=\small\bfseries, align=center},
    arr/.style={-Stealth, thick, deepblue},
    arrno/.style={-Stealth, thick, dangerred},
]

\node[term, fill=inputyellow, draw=deepblue] (start) at (0,7.0) {输入 障碍物集合};

\node[box] (init) at (0,5.2) {初始化路径边界为车道边};

\node[dec] (loop) at (0,3.2) {遍历障碍物 $O_i$};

\node[box] (decide) at (-5,1.2) {判断避让方向\\ $d_i\in\{L,R\}$};

\node[boximp] (update) at (0,1.2) {更新 $l^+$ 或 $l^-$};

\node[dec] (check) at (0,-1.2) {通行带宽 $W(s) < W_{\text{ego}}$?};

\node[term, fill=dangerred!40, draw=dangerred] (blocked) at (5.5,-1.2) {BLOCKED};
\node[term] (ok) at (0,-3.4) {输出 path\_boundary};

% 流
\draw[arr] (start) -- (init);
\draw[arr] (init) -- (loop);
\draw[arr] (loop) -- (decide);
\draw[arr] (decide) -- (update);
\draw[arr] (update) -- (check);
\draw[arrno] (check) -- node[above, font=\footnotesize, dangerred] {是} (blocked);
\draw[arr] (check) -- node[left, font=\footnotesize, deepblue] {否} (ok);

% 反馈回 loop（遍历下一个）
\draw[arr] (update.east) -- ++(1.8,0) |- (loop.east);
\node[font=\scriptsize, deepblue] at (2.6,2.2) {下一个 $O_{i+1}$};

% 侧注
\node[font=\footnotesize, align=left, anchor=north west, deeppink!80!black] at (3.8,4.5)
  {核心问题：决策逐障碍物独立做出，\\后续障碍物无法修正前序决策。};

\end{tikzpicture}
}
\caption{PathBoundsDecider相关边界构建流程}
\label{fig:path_bounds_flow}
\end{figure}

上述流程有两个关键特征。第一，避让方向决策是逐障碍物独立进行的，每个障碍物仅根据自身位置相对于参考线的横向偏移做出左绕或右绕判定，不考虑其他障碍物已做出的决策对剩余通行空间的影响。第二，障碍物过滤函数\texttt{IsWithinPathDeciderScopeObstacle}将所有动态障碍物排除在路径边界决策之外，仅由后续的速度规划阶段处理。这两个特征在多障碍物场景下会引发第\ref{chap:framework}章所分析的问题。

\section{SpeedBoundsDecider概述}

\texttt{SpeedBoundsDecider}在路径已确定的前提下将障碍物映射到ST图中，构建ST边界作为速度QP的约束。其处理对象是路径规划阶段输出的\texttt{PathData}，遍历路径上的离散点逐一检查与障碍物的碰撞关系。本节先梳理其构建流程，再将其与\texttt{PathBoundsDecider}相关边界构建逻辑在架构层面进行对比。

\subsection{ST边界构建流程}

\texttt{SpeedBoundsDecider}的处理流程分为三步。第一步，沿已确定的路径逐点检查与所有障碍物的碰撞关系，将每个障碍物在ST图中的占据区域表示为一个多边形。静态障碍物在ST图中表现为沿时间轴延伸的矩形带状区域，动态障碍物表现为斜向的平行四边形，其斜率由障碍物沿参考线方向的速度决定。第二步，在每个时间步$t_j$处，汇总所有障碍物ST多边形的约束，综合YIELD、FOLLOW和OVERTAKE三类决策，分别给出纵向位置的上界$s_j^{ub}$和下界$s_j^{lb}$。YIELD和FOLLOW决策要求自车位于障碍物之后，约束体现为$s_j^{ub}$的收紧；OVERTAKE决策要求自车位于障碍物之前，约束体现为$s_j^{lb}$的抬升。第三步，在ST边界之上叠加纵向安全距离$d_s$，最终输出供后续速度优化器使用的$[s_j^{lb}, s_j^{ub}]$序列。

\subsection{与路径边界构建逻辑的架构差异}

\texttt{SpeedBoundsDecider}与\texttt{PathBoundsDecider}相关边界构建逻辑分处路径-速度解耦架构的不同阶段，两者有三处主要差异。

第一，处理的障碍物范围不同。\texttt{PathBoundsDecider}相关边界构建逻辑仅处理静态障碍物，而\texttt{SpeedBoundsDecider}同时处理静态和动态障碍物。动态障碍物在路径阶段被完全排除，只能在速度阶段通过ST图处理。

第二，决策协调机制不同。\texttt{PathBoundsDecider}相关边界构建逻辑对每个障碍物独立做出左绕/右绕的贪心决策，没有全局协调。\texttt{SpeedBoundsDecider}之后的\texttt{PathTimeHeuristicOptimizer}即DP启发式搜索，在ST图中对YIELD/OVERTAKE的组合进行全局搜索，将DP输出作为后续QP的暖启动参考，避免了逐障碍物贪心策略在速度层继续累积。

第三，信息传递方向单向。路径阶段的输出即\texttt{PathData}传递给速度阶段作为已知条件，但路径阶段做出决策时的上下文信息，例如\emph{该路径依赖某动态障碍物在特定时刻移开}这一时间有效性前提，并未传递给速度阶段。速度阶段仅看到路径的几何形态，无法获知其背后的决策依据。路径决策和速度规划之间因此缺乏时空约束的协调，这是第\ref{chap:st_corridor}章所分析的ST层问题的架构根源。第\ref{chap:st_corridor}章提出的时变通行带框架输出时间有效性约束，以ST上界的形式与\texttt{SpeedBoundsDecider}的输出在QP层面融合，补上路径-速度之间缺失的约束传递。

\section{本章小结}

本章从源码层面梳理了Apollo Planning模块的主要组件。\texttt{PathBoundsDecider}相关边界构建逻辑在Apollo 11.0中主要由\texttt{PathBoundsDeciderUtil}承载，其特征体现在三点，逐障碍物独立决定避让方向、对所有障碍物使用统一安全裕度、将动态障碍物完全排除在路径决策之外。\texttt{SpeedBoundsDecider}借助DP启发式搜索在ST图上完成纵向组合决策的全局协调，但从路径阶段到速度阶段的信息传递仅是单向的，路径决策所依赖的时间有效性前提无法随\texttt{PathData}一同下传。这些特征在多障碍物密集场景下会集中表现为路径决策的保守与协同缺位，构成后续章节问题识别的源码依据。
